\uuid{0MXI}
\titre{Division euclidienne}
\auteur{bodin ; hanani ; mzari}
\organisation{exo7}
\licence{CC BY-NC-SA 4.0}
\datecreate{2020-01-20}
\niveau{L1}
\module{Algèbre}
\chapitre{Polynôme, fraction rationnelle}
\sousChapitre{Division euclidienne}
\difficulte{1}
\theme{Arithmétique des polynômes, division euclidienne}
\exo7qcmid{lille-1.polynomes.009}

\contenu{
\begin{question}
\qclassification{105.01,105.02}
\qtags{motcle=division euclidienne}
Soient $A,B$ deux polynômes, avec $B$ non nul.
Soit $A = B \times Q + R$ la division euclidienne de $A$ par $B$.
\begin{answers}
    \good{Un tel $Q$ existe toujours.}

    \good{S'il existe, $Q$ est unique.}

    \good{On a toujours $\deg Q \le \deg A$.}

    \bad{On a toujours $\deg Q \le \deg B$.}
\end{answers}
\begin{explanations}
La division euclidienne $A = B \times Q + R$ existe toujours, $Q$ et $R$ sont uniques et bien sûr $\deg Q \le \deg A$.
\end{explanations}
\end{question}
}
